\section{Related Work}
\label{sec:related}

\subsection{Cross-Domain Industrial Anomaly Detection}
Deep transfer learning for industrial time-series anomaly detection is surveyed in \cite{survey2023}, which organizes methods by transfer setting and lists data gaps as a main obstacle. Two recent systems define the state of the art in feature-space transfer: CICADA \cite{cicada2025} learns interpretable mixture-of-experts codes that expand to unseen target domains, and DACAD \cite{dacad2024} combines anomaly injection with contrastive unsupervised domain adaptation. Both optimize detection scores; neither produces calibrated decision thresholds, and both assume aligned feature spaces that ignore frequency structure. Our study takes their score-producing pipeline as given and asks the complementary reliability question: at what threshold should the score be deployed, and does the answer survive a domain change?

\subsection{Time-Series Foundation Models Under Shift}
Chronos \cite{chronos2024} and PatchTST-style backbones \cite{patchtst2023} popularized reusable sequence representations for forecasting, and they are frequently used zero-shot for anomaly scoring. Two recent analyses temper this: \cite{freqmatters2025} isolates \emph{spectral shift}---a mismatch between downstream dominant frequencies and pretraining spectra---as a driver of foundation-model failure, and \cite{foundational2025} argues that the diversity of real series undermines the foundation-model premise. Neither proposes a deployable thresholding fix. Our benchmark deliberately keeps native sampling rates so that spectral shift is preserved rather than resampled away.

\subsection{Conformal Prediction for Time Series}
Split conformal prediction \cite{lei2018} gives distribution-free marginal coverage from exchangeable calibration data; time series and domain shift violate exchangeability. Adaptive conformal inference (ACI) \cite{aci2021} replaces exchangeability with an online feedback rule on the conformal quantile and guarantees long-run miscoverage under arbitrary dependence, and \cite{dtaci2022} systematizes update variants (including the DtACI multiplicative-plus-additive rule) for temporal streams. Conformal PID control \cite{conformalpid2023} casts the update as a controller and adds proportional--derivative terms. Extreme-value thresholding \cite{siffer2017} is the classical streaming alternative, controlling tail exceedance with a generalized-Pareto fit. All of this work studies temporal drift within a domain; cross-domain (cross-population) calibration on real industrial data, to our knowledge, has not been evaluated systematically. We provide that evaluation, including the cases where the conformal layer succeeds and the cases where it cannot help.

\subsection{Parameter-Efficient Adaptation}
Test-time adaptation and parameter-efficient fine-tuning---e.g., the robust frequency-domain adapter of PETSA \cite{petsa2025}---are natural candidates for closing cross-domain gaps. In our earlier ablations, test-time adaptation and gating brought no measurable detection gain at this scale, and we therefore keep the encoder fixed and study the calibration layer in isolation; adaptation methods remain compatible with the framework as score providers.
